\input{../tex-inputs/latex-leseansicht-vorspann}

\section[Carl von Torresani an Arthur Schnitzler, {[}zwischen 13. 5. und 31. 12. 1891?{]}]{L04323 Carl von Torresani an Arthur Schnitzler, {[}zwischen 13. 5. und 31. 12. 1891?{]}}
\nopagebreak\mylabel{L04323v}
\rehead{ }\normalsize\beginnumbering\briefempfaengerindex{Schnitzler, Arthur@\textsc{Schnitzler, Arthur}!zzzTorresani-Lanzenfeld, Carl von@\emph{von Carl von Torresani-Lanzenfeld}!1891-12-311@{{[}zwischen 13. 5. und 31. 12. 1891?{]}}|(be}
\toendnotes[C]{\smallbreak\pagebreak[2]}
\correspDesc{Versand  durch Carl von Torresani im Zeitraum [zwischen 13. 5. und 31. 12. 1891?] in Venedig
\newline{}Erhalt  durch Arthur Schnitzler in Wien}\toendnotes[C]{\smallbreak}
\Standort{CUL, Schnitzler, B 106.}
\physDesc{Postkarte, 589 Zeichen
\newline{}Handschrift: schwarze Tinte, deutsche Kurrent
\newline{}Versand: Stempel: »\nobreak{}\oindex{Venedig@\textbf{Venedig}|pwk}V\textcolor{gray}{enezia}, 5–6\nobreak{}«.  }\toendnotes[C]{\smallbreak}\pstart{}{\pb}\textcolor{gray}{\textbf{A}}n Herrn D\textsuperscript{r} Arthur
                  Schnitzler\pend{}\pstart{}Wohlgeb\pend{}\pstart{}\textsc{Wien\oindex{Wien@\textbf{Wien}, \emph{Verwaltungsgebiet}|pw}}\pend{}\pstart{}I Giſelaſtraße 11\oindex{Wien@\textbf{Wien}!I., Innere Stadt@\textbf{I., Innere Stadt}!Kärntnerring 12/Bösendorferstraße 11@\textbf{Kärntnerring 12/Bösendorferstraße 11}, \emph{Wohngebäude}|pw}\pend{}\pstart{}\textsc{Austria\oindex{Österreich@\textbf{Österreich}|pw}}\pend{}{\bigskip}\vspace{1em}
\pstart
           \noindent{}{\pb}Ich habe, geſchätzteſter Herr
                  Doktor, meine Zuſage nicht vergeſſen; theile Ihnen demzufolge meine
               Adreſſe mit: \textsc{\uline{Venezia, Riva dei Schiavoni, nuova Casa Kirsch}\oindex{Albergo Metropole@\textbf{Albergo Metropole}, \emph{Hotel}|pw}}; gleichzeitig geht auch \label{K_L04323-1v}\edtext{das verſprochne Buch\pwindex{Torresani-Lanzenfeld, Carl von 19.\,4.\,1846 Mailand – 16.\,4.\,1907 Nago-Torbole@\textsc{Torresani-Lanzenfeld, Carl von} (19.\,4.\,1846 Mailand – 16.\,4.\,1907 Nago-Torbole), \emph{Schriftsteller, Offizier}!Juckercomtesse. Roman aus der Gesellschaft@\strich\emph{Die Juckercomtesse. Roman aus der Gesellschaft}|pwuv}}{\lemma{\textnormal{\emph{das versprochne Buch}}}\Cendnote{\textnormal{Vermutlich handelt es sich um Torresanis\pwindex{Torresani-Lanzenfeld, Carl von 19.\,4.\,1846 Mailand – 16.\,4.\,1907 Nago-Torbole@\textsc{Torresani-Lanzenfeld, Carl von} (19.\,4.\,1846 Mailand – 16.\,4.\,1907 Nago-Torbole), \emph{Schriftsteller, Offizier}|pwk} Roman
                     \emph{Die Juckercomtesse}\pwindex{Torresani-Lanzenfeld, Carl von 19.\,4.\,1846 Mailand – 16.\,4.\,1907 Nago-Torbole@\textsc{Torresani-Lanzenfeld, Carl von} (19.\,4.\,1846 Mailand – 16.\,4.\,1907 Nago-Torbole), \emph{Schriftsteller, Offizier}!Juckercomtesse. Roman aus der Gesellschaft@\strich\emph{Die Juckercomtesse. Roman aus der Gesellschaft}|pwk}, dessen Erscheinen das \emph{Börsenblatt für den Deutschen  Buchhandel}\pwindex{Börsenblatt für den Deutschen Buchhandel@\emph{Börsenblatt für den Deutschen Buchhandel}|pwk} am 25. 4. 1891 vermeldete.}}}\label{K_L04323-1} an Sie ab. – Bisher habe ich von meiner
               Reiſe nicht viel Freude gehabt; regneriſches u kaltes Wetter, u mancherlei
               Widerwärtigkeiten anderer Art. Hoffentlich war dies mein Polykrates-Ring\pwindex{\textcolor{red}{\textsuperscript{XXXX indx1}}!Ring des Polykrates@\strich\emph{Der Ring des Polykrates}|pwv}, u \label{K_L04323-2v}\edtext{kein Fisch}{\lemma{\textnormal{\emph{kein Fisch}}}\Cendnote{\textnormal{In Schillers Ballade \emph{Der
                     Ring des Polykrates}\pwindex{\textcolor{red}{\textsuperscript{XXXX indx1}}!Ring des Polykrates@\strich\emph{Der Ring des Polykrates}|pwk} wirft der überaus erfolgreiche und glückliche
                  Protagonist sein liebstes Schmuckstück, einen Ring, ins Meer, um zu verhindern,
                  dass die Götter missgünstig wegen seines anhaltenden Glücks werden und sich gegen
                  ihn wenden. Als am Folgetag sein Koch mitteilt, der weggeworfene Ring habe sich
                  zufällig in einem gefangenen Fisch wiedergefunden, ist klar, dass
                  Polykrates Untergang nach dieser erneuten Begünstigung durch das Schicksal
                  nicht aufgehalten werden kann.}}}\label{K_L04323-2} bringt mir ihn wieder.\pend
           \pstart Beſte Grüße \label{K_L04323-3v}\edtext{allen Herrn}{\lemma{\textnormal{\emph{allen Herrn}}}\Cendnote{\textnormal{Das Korrespondenzstück ist nicht datiert.
                  Schnitzler erwähnte Torresani\pwindex{Torresani-Lanzenfeld, Carl von 19.\,4.\,1846 Mailand – 16.\,4.\,1907 Nago-Torbole@\textsc{Torresani-Lanzenfeld, Carl von} (19.\,4.\,1846 Mailand – 16.\,4.\,1907 Nago-Torbole), \emph{Schriftsteller, Offizier}|pwk} im \emph{Tagebuch}\pwindex{Schnitzler, Arthur 15.\,5.\,1862 Wien – 21.\,10.\,1931 ebd.@\textsc{Schnitzler, Arthur} (15.\,5.\,1862 Wien – 21.\,10.\,1931 ebd.), \emph{Schriftsteller, Mediziner}!Tagebuch@\strich\emph{Tagebuch}|pwk} zum ersten Mal am 29. 3. 1891 und
                  vertiefte die Bekanntschaft offenbar am 12. 5. 1891. Vermutlich wurde die vorliegende Karte bald im Anschluss
                  verfasst, denn die auszurichtenden Grüße gehen mit Felix Salten\pwindex{Salten, Felix 6.\,9.\,1869 Budapest – 8.\,10.\,1945 Zürich@\textsc{Salten, Felix} (6.\,9.\,1869 Budapest – 8.\,10.\,1945 Zürich), \emph{Schriftsteller, Journalist, Chefredakteur}|pwk}, Eduard
                     Michael Kafka\pwindex{Kafka, Eduard Michael 11.\,3.\,1869 Wien – 6.\,8.\,1893 Brünn@\textsc{Kafka, Eduard Michael} (11.\,3.\,1869 Wien – 6.\,8.\,1893 Brünn), \emph{Redakteur}|pwk} und Jaques Joachim\pwindex{Joachim, Jaques 24.\,11.\,1866 Wien – 7.\,11.\,1925 ebd.@\textsc{Joachim, Jaques} (24.\,11.\,1866 Wien – 7.\,11.\,1925 ebd.), \emph{Rechtswissenschaftler, Rechtsanwalt, Herausgeber}|pwk} an
                  einen Zirkel, der im Jahr 1891 um die Zeitschrift \emph{Moderne Rundschau}\pwindex{Moderne Dichtung. Monatsschrift für Literatur und Kritik@\emph{Moderne Dichtung. Monatsschrift für Literatur und Kritik}|pwk} herum bestand,
                  deren Erscheinen bereits im Dezember 1891 eingestellt wurde.
                  (Begegnungen mit Joachim\pwindex{Joachim, Jaques 24.\,11.\,1866 Wien – 7.\,11.\,1925 ebd.@\textsc{Joachim, Jaques} (24.\,11.\,1866 Wien – 7.\,11.\,1925 ebd.), \emph{Rechtswissenschaftler, Rechtsanwalt, Herausgeber}|pwk} erwähnte Schnitzler schon nach dem 25. 6. 1891 keine
                  mehr.)}}}\label{K_L04323-3}, insbeſ. \textsc{Salten\pwindex{Salten, Felix 6.\,9.\,1869 Budapest – 8.\,10.\,1945 Zürich@\textsc{Salten, Felix} (6.\,9.\,1869 Budapest – 8.\,10.\,1945 Zürich), \emph{Schriftsteller, Journalist, Chefredakteur}|pw}}, \textsc{Kafka\pwindex{Kafka, Eduard Michael 11.\,3.\,1869 Wien – 6.\,8.\,1893 Brünn@\textsc{Kafka, Eduard Michael} (11.\,3.\,1869 Wien – 6.\,8.\,1893 Brünn), \emph{Redakteur}|pw}}, \textsc{Joachim\pwindex{Joachim, Jaques 24.\,11.\,1866 Wien – 7.\,11.\,1925 ebd.@\textsc{Joachim, Jaques} (24.\,11.\,1866 Wien – 7.\,11.\,1925 ebd.), \emph{Rechtswissenschaftler, Rechtsanwalt, Herausgeber}|pw}}; u nehmen Sie selbſt die herzlichſten Grüße von \spacefill\mbox{Torresani}\pend{}\selectlanguage{ngerman}\endnumbering\briefempfaengerindex{Schnitzler, Arthur@\textsc{Schnitzler, Arthur}!zzzTorresani-Lanzenfeld, Carl von@\emph{von Carl von Torresani-Lanzenfeld}!1891-05-131@{{[}zwischen 13. 5. und 31. 12. 1891?{]}}|)be}\mylabel{L04323h}
\begin{anhang}
\end{anhang}\newcommand{\dateiname}{L04323}\newcommand{\titel}{Carl von Torresani an Arthur Schnitzler, [zwischen 13. 5. und 31. 12. 1891?]}\newcommand{\editorInnen}{Herausgegeben von Jahnke, SelmaMüller, Martin Anton}\input{../tex-inputs/latex-leseansicht-abspann}
